% !TeX root = ../papers/nc-p3.tex
\section{節点または尖点の曲線上の局所印と長さ六の場合}
\label{sec:finite-nodal-cuspidal-marks}
\begin{theorem}\label{thm:finite-nodal-cuspidal-marks}
有限な共通切断で印の接次元が2であり、その局所環が通常の節点または通常の尖点を持つ
平面曲線上の非零因子による商なら、長さ3の階数一切断では常に偶側の反対対称性が成立する。
反対階数8の有限窓では、長さ6の局所因子が一つでもあれば全次数の正則性が有限に決まる。
従って残余の有限窓の閉点数は高々二つである。
\end{theorem}
\begin{proof}[曲線の導手による印での評価]
印の局所環を $P=C/(g)$ とする。完備な節点は $C=k[[u,v]]/(uv)$ で、
正規化の二枝での $g$ の位数を $n_j$、各枝の極大イデアルの最小位数を $d_j=1$ とする。
尖点は $C=k[[t^2,t^3]]$、$n=\operatorname{ord}_t g$、$d=2$ とする。
いずれも正規化で各位数が $d_j$ 以上の元は $C$ に入る。
従ってある持上げの各位数が $n_j+d_j$ 以上なら、それは $gC$ に入る。
位置1,3の $l,r$ の各列は $P$ の極大イデアルを生成する。$g$ は極大イデアルの二乗に入り、
各枝で両列の最小位数がちょうど $d_j$ である。

$Z$ を印の二行列 $A,B$ の規格化した差とする。各枝で $Z$ の最小位数を $s$、一様化変数を $t$ と書く。
$s<n_j-d_j$ なら、法 $t^{n_j-s-d_j}$ で $A,Z/t^s$ は独立な二元であり、
$l/t^{d_j},r/t^{d_j}$ の剰余で非零な二方向を消す。共通の独立性は $Z$ を固定超平面に正規化したことによる。
従ってこれらは相対割線の基底である。Gorenstein三重点の模型では三方向のスキームで消える二次式 $Q$ に対し、
$Q(l/t^{d_j})=0\bmod t^{n_j-s-d_j}$ である。ゆえに $Q(l)Z$ の各位数は
$2d_j+s+(n_j-s-d_j)=n_j+d_j$ 以上となる。
$s\ge n_j-d_j$ なら同じ位数は二次式の次数だけで足りる。導手により元の $P$ 上で $Q(l)Z=0$ である。
定理\ref{thm:finite-marked-artin-secants}のイデアルの列と完全跡の議論から対称性が従う。

非Gorenstein三重点では定理\ref{thm:finite-marked-planar-triple}の例外平面の式を $T,l_0$ と書く。
上の相対割線では $T(A)$ が零で、正規化した左方向の $l_0$ も零である。
第二辺関係の評価 $f$ は $Z$ の成分の局所線型結合であるから、
$T f,l_0 f$ の枝での位数は少なくとも $n_j-d_j$ である。$s\ge n_j-d_j$ でも同じである。
従って $\mathfrak m^2 T f=\mathfrak m^2 l_0 f=0$ が導手で従う。印以外でこれらは零であり、
各々の支持環はGorensteinで極大イデアルの二乗が零なので長さ2以下である。
隣の一次座標の全射と跡直交で両方とも零になる。よって既出の長さ6の完全交叉 $Y$ 上の誤差議論が適用できる。

そこで得る誤差イデアルの一生成元を再び $f$ とする。
印で $f\ne0$ なら $P/\operatorname{ann}f$ は曲線状であることを一般に確かめる。
$Z=fV'$ とし $V'$ の剰余は $A$ の剰余 $M$ と独立である。
この商で $V'$ は $l,r$ の両方を消す。もし接次元2なら、両者の一次方向はそれぞれ
$M$ の左・右核の二平面全体を張る。$V'$ の剰余もこれら両平面を消し、階数1の $M$ の倍となって矛盾する。
よってこの商で一様化変数を割り、一つだけ局所長さを落として相対割線を得る。
$Y$ の長さ5以上での余因子の消滅は既出の証明と同じ矛盾を与える。
\end{proof}

\begin{proof}[長さ六の平面局所環]
前定理で扱われない長さ6の平面Gorenstein局所環は Hilbert 関数 $(1,2,2,1)$ を持つ。
実際二次部分が三次元ならそれがsocleとなってしまう。
有限辺は六次元滑らかな多様体での六双一次式の完全交叉なので、
完備局所環で独立な一次式四つを消せば $P=k[[x,y]]/(F,G)$ であり、両式の位数は2以上である。
二次初期式は一次元を張る。$F$ に非零の二次初期式を持たせ、$G$ の位数は3以上とする。
$F$ の二次が二つの異なる方向を持てば通常の節点である（異なる一次因子を逐次持ち上げる）。
二次が $y^2$ なら、$F$ に $G$ の適当な定数倍を足して通常の尖点にできる。
実際両式の $x^3$ 係数が零と仮定する。Weierstrassの準備と除法で
\[
                F=y^2+a(x)y+b(x),\qquad G=c(x)+d(x)y
\]
とでき、各位数は $2,4,4,2$ 以上である。長さは自由な二次の $k[[x]]$ 加群での
$G$ の乗法行列の行列式 $c^2-acd+bd^2$ の位数なので8以上となり矛盾する。
非零の $x^3$ 項があれば平方完成と単元の三乗根で通常の尖点となる。
これで長さ3の切断は最初の証明で尽きた。

長さ4の場合は定理\ref{thm:finite-stretched-planar-marks}の二次小行列式の議論を使う。
$n=3$ または局所 $(1,2,1)$ 型の $n=4$ では、非零であるはずの隣接語 $f=H\chi(a)$ は
印以外で $(\eta)$ に入り、印で $\mathfrak m^2 f\subset(\eta)$ である。
印が $Z$ の台から外れれば印の寄与は高々2であり、先の証明の通り矛盾する。
印が台に入るとき $n=4$ のこの場合は $Z$ が印に集中し、$f$ は非単元なので
支持長は6未満である。$n=3$ では $\operatorname{length}(f)\le3+3=6$ であり、
等号ならその支持は $Z$ を含む。長さ6の隣接語の偶側の包は平面なので、
$Z$ を張る平面の式 $\chi(a)$ は $f$ を消す。この式が印で一次の非零方向を持てば
$f$ の局所支持環は曲線状であり、$\mathfrak m^2 f\subset(\eta)$ と合わせ全長は5以下となる。
そうでなければ $\chi(a)\in\mathfrak m^2$ で $f$ の局所寄与は2以下である
（$\mathfrak m^4=0$ ゆえ $f$ の局所支持環の極大イデアルの二乗は零）。
いずれも矛盾する。最後に $n=4$ で曲線状因子の積の場合、四つの独立な偶三次式
$\operatorname{adj}_B(a)$ は印以外で零であり、印では乗法階数1より $\mathfrak m^2$ に入る。
その長さはちょうど3なので矛盾する。共線因子の排除と短い切断の場合は既出の通りである。
長さ20で三局所因子があれば一つは長さ6以下であるから最後の主張も従う。
\end{proof}

\begin{corollary}\label{cor:finite-no-four-on-line}
四辺が有限な標準有限Frobenius窓の各位置の点スキームは、
どの射影直線とも長さ3以下で交わる。反対階数の制限はない。
\end{corollary}
\begin{proof}
直線上に長さ4の因子があれば、その前または後の座標の包を $\mathbb P^{a-1}$ とする。
隣接辺の六式をこの包と直線の積に制限し、一般の $a$ 本を取れば
有限基点なので固有交叉となり、その長さは $a$ である。従って $a\ge4$ であり、両隣の座標は
因子の全長さ4の環をそれぞれ張る。定理\ref{thm:finite-artin-kernel-symmetry}の巡回行列の証明によって
純関係ができ、有限性に反する。
\end{proof}
